\Needspace{12\baselineskip}
\originalchapter{映射、维数与混沌尺度}
\originalsection{Hénon映射与局部方向}
二维离散映射 $x_{n+1}=1+y_n-\alpha x_n^2$、$y_{n+1}=\beta x_n$ 的Jacobian行列式为 $-\beta$。$|\beta|<1$ 时面积每步按固定比例收缩；折叠来自二次项。固定点满足 $\alpha x^2+(1-\beta)x-1=0$，稳定性由Jacobian特征值的模判定，而不是实部。

映射附近的伸长和收缩方向可用特征向量求出。吸引子在鞍固定点附近常沿其不稳定流形展开，尽管全局长期轨迹仍被吸引到一个有限集合。将“固定点不稳定”误读为“所有轨迹都逃向无穷远”是不正确的。

\originalsection{倍周期与普适范围}
Logistic映射 $x_{n+1}=r x_n(1-x_n)$，$0\le r\le4$ 保持 $[0,1]$ 不变。非零固定点 $1-1/r$ 的乘子 $2-r$ 在 $r=3$ 越过 $-1$，产生倍周期。周期 $p$ 的乘子是沿周期 $\prod_{j=0}^{p-1}f'(x_j)$，有一点位于二次极值 $x=1/2$ 时乘子为零，称超稳定周期。

沿二次单峰映射的倍周期序列，参数间隔比可趋近Feigenbaum常数约4.669。极值为四次而非二次、非光滑或多峰系统未必属于同一普适类。有限级的间隔比只显示收敛趋势，不能代替极限存在和普适性的证明。所选PS10.1实际计算到 $2^7$ 周期。

间歇路线常在鞍结附近经过很慢的层流阶段。例如 $\dot x=\epsilon+x^2$ 穿越瓶颈的时间 $\int\dd x/(\epsilon+x^2)\sim\pi/\sqrt\epsilon$，之后快速回返可形成间歇爆发。准周期破裂则涉及多个相位，与倍周期机制不同，不能给所有混沌图同贴一种路线标签。

\originalsection{分形和Lyapunov指数}
自相似集合每轮有 $m$ 个缩小为原长 $r$ 的相似片，维数由 $mr^D=1$ 得 $D=\ln m/\ln(1/r)$，需要各片分离等条件。箱维数、相关维数和Lyapunov维数的定义不同，不能无条件互换。

对于流的切向扰动 $\dot{\delta\mathbf x}=J\delta\mathbf x$，最大Lyapunov指数 $\lambda_{max}=\lim_{t\to\infty}t^{-1}\ln(\|\delta\mathbf x(t)\|/\|\delta\mathbf x(0)\|)$ 需不断将扰动维持在线性范围。直接比较两条很近轨迹，差异最终饱和，后期对数斜率趋零并不表示没有混沌。

相关维数以 $C(\epsilon)\sim\epsilon^{D_2}$ 的小尺度斜率估计。必须排除时间相邻点形成的平滑轨迹段，且有限噪声、有限点数会造成假平台。真正的吸引子维数不能由一次任意尺度区间拟合确定，本版未把PS9.3全参数维数估计列为已完成题。

\originalsection{H\'enon的伸展、折叠与面积收缩}
映射可分为二次弯曲、横向缩放和坐标交换，构成伸展后折返的机制。固定点由二次式给两根，Jacobian $J=\begin{pmatrix}-2\alpha x&1\\\beta&0\end{pmatrix}$，乘子满足 $\lambda^2+2\alpha x\lambda-\beta=0$。$|\beta|<1$ 的面积收缩不妨碍其中一个 $|\lambda|>1$。

对正根固定点，乘子 $-1$ 条件给 $\alpha_c=3(1-\beta)^2/4$（题面参数域内），产生翻转失稳。$\alpha=1.4,\beta=.3$ 时不稳定特征方向斜率由 $dy/dx=\beta/\lambda_u$ 给，实际数值约 $-.155946$。曲线附近点云的全局主轴不必等于局部流形切线；应从固定点微小扰动和切向迭代比较。

\originalsection{分形维数的定义和实际尺度窗口}
Cantor集每次保留两段、各长原来的 $1/3$，第 $n$ 层需要 $2^n$ 箱、箱宽 $3^{-n}$，所以 $D=\ln2/\ln3$。面积/长度为零并不表示点集有限；维数衡量覆盖数量随分辨率增长，不是将几何图看成“介于一维和二维”就有数值证明。

箱维数用 $N(\epsilon)$，相关维数用随机点对近距离概率 $C(\epsilon)$。若点权重均匀且规则自相似可相同，非均匀测度下通常 $D_2\le D_0$。数值要在大尺度有限集范围和小尺度噪声/点数稀疏之间寻找稳定窗口；排除时间近邻的Theiler窗避免把一条光滑轨迹的短片段误当吸引子统计。

\originalsection{Lyapunov谱的切向积与QR}
映射切向扰动 $\delta x_n=J_{n-1}\cdots J_0\delta x_0$。单个矩阵的瞬时特征值不能代替乘积长时增长，尤其不同 $J_n$ 不对易。为求完整谱，传播一组正交扰动，按固定间隔QR分解，累计 $\ln|R_{ii}|$ 并除总时间；不断正交化避免所有列都沿最大方向对齐。

流的完整谱之和等于沿轨道平均散度，在有界非固定点自治轨道通常有一个零指数对应时间平移。对 $d$ 维吸引子，Kaplan--Yorke公式 $D_{KY}=j+\sum_{i=1}^j\lambda_i/|\lambda_{j+1}|$ 是Lyapunov维数，和相关维数的相等需要附加理论条件。两个指数的和可检查H\'enon面积因子 $\lambda_1+\lambda_2=\ln|\beta|$。

\originalsection{Logistic二周期与超稳定序列}
非零固定点在 $r=3$ 翻转。去掉 $f(x)=x$ 因子后，$f^2(x)=x$ 给二周期
\[
x_\pm=\frac{r+1\pm\sqrt{(r-3)(r+1)}}{2r}.
\]
周期乘子为 $f'(x_+)f'(x_-)=-r^2+2r+4$，模小于1给 $3<r<1+\sqrt6$。这是从固定点失稳到二周期再失稳的完整一层推导，不只观察分岔图。

另一约定 $f(x)=4\mu x(1-x)$ 用 $\mu=r/4$。超稳定 $2^n$ 周期要求 $f^{2^n}(1/2)=1/2$ 并排除低周期根；实际程序以逐级区间定位找到 $n=1$ 到7的主序列，不能把所有方程根中第一个随便称为下一代。所选原题给全数值和间隔比。

\originalsection{倍周期重整化的选择结构}
在二次极值附近，迭代两次再缩放坐标使峰形归一，定义重整化算子 $\mathcal R[f](x)=-\alpha f(f(-x/\alpha))$（这里明确采用这一符号约定）。临界普适函数满足 $g=\mathcal R[g]$，$g(0)=1$、极值二次。线性化 $\mathcal R[g+\epsilon h]=g+\epsilon\mathcal Lh+\cdots$，最相关本征值 $\delta$ 控制参数距离每倍周期缩小的比例；坐标缩放 $\alpha$ 控制空间尺度。

这解释为什么参数间隔比和几何尺寸比是两个不同常数。数值的有限级比趋近约4.669不证明普适函数存在，但与理论机制一致。高阶极值、非光滑或多峰改变重整化类别，不应把同一常数机械应用。

\Needspace{12\baselineskip}
\originalchapter{间歇和准周期的混沌路线}
\originalsection{慢瓶颈、锁频与环面破裂}
Type-I间歇的回返映射近相切为 $x_{n+1}=x_n+\epsilon+x_n^2$。小 $\epsilon>0$ 时每步慢漂，连续近似通过瓶颈耗步数 $\int\dd x/(\epsilon+x^2)\sim\pi/\sqrt\epsilon$，之后全局重新注入产生快爆发。慢阶段时长的分布还依赖重新注入机制，因此平均尺度与完整分布不是同一个结论。

准周期在二维环面上有两个相位，旋转数为频率比；有理比对应锁频，不成有理比轨道稠密。圆映射 $\theta_{n+1}=\theta_n+\Omega-(K/(2\pi))\sin(2\pi\theta_n)\pmod1$ 展示锁频舌与准周期，强耦合可破坏环面。Ruelle--Takens路线涉及高维耗散系统中环面的失稳，不能按“多出三个谱峰”就证明其前提已满足。
